\subsection{Bounded measures}

Let X be a locally compact space. Since the topology on \(\mathcal{K}(X;\mathbf{C})\) induced by that of \(\mathcal{C}^{b}(X;\mathbf{C})\) is coarser than the direct limit topology on \(\mathcal{K}(X;\mathbf{C})\) , a measure on X is not necessarily continuous for the topology of uniform convergence in X.

DEFINITION 3. --- A measure on a locally compact space X is said to be bounded if it is continuous on \(\mathcal{K}(\mathrm{X};\mathbf{C})\) for the topology of uniform convergence.

It comes to the same to say that there exists a finite number \(M \geqslant 0\) such that, for every function \(f \in \mathcal{K}(\mathbf{X}; \mathbf{C})\) ,

\[
| \mu (f) | \leqslant \mathrm{M} \| f \|\tag{19}
\]

(where \(\| f\|\) is defined by formula (3) of No.~2).

To say that \(\mu\) is a bounded measure thus signifies that \(\mu\) belongs to the dual of the space \(\mathcal{K}(\mathrm{X};\mathbf{C})\) normed by \(\| f\|\) ; we shall denote this dual by \(\mathcal{M}^1 (\mathrm{X};\mathbf{C})\) (or simply \(\mathcal{M}^1 (\mathrm{X})\) when no confusion can result). We know that \(\mathcal{M}^1 (\mathrm{X};\mathbf{C})\) is equipped with a norm, \(\| \mu \|\) being the smallest of the numbers \(M\geqslant 0\) for which the inequality (19) holds for every function \(f\in \mathcal{K}(\mathrm{X};\mathbf{C})\) , or again,

\[
\| \mu \| = \sup _ {\| f \| \leqslant 1, f \in \mathcal{K} (\mathrm{X}; \mathbf {C})} | \mu (f) |.\tag{20}
\]

Equipped with this norm, \(\mathcal{M}^{1}(X;\mathbf{C})\) is known to be a Banach space (TVS, III, §3, No.~8, Cor. 2 of Prop. 12).

The definition of \(\|\mu\|\) by the formula (20) may be extended to every measure \(\mu\) on X and, by an abuse of language, \(\|\mu\|\) is again said to be the norm of \(\mu\) ; for \(\mu\) to be bounded, it is necessary and sufficient that \(\|\mu\|\) be finite.

If X is compact, then every measure on X is bounded.

Examples. --- 1) The measure \(\varepsilon_{a}\) defined by a unit mass at a point \(a \in X\) is bounded, and \(\| \varepsilon_{a} \| = 1\) .

\begin{enumerate}
\def\labelenumi{\arabic{enumi})}
\setcounter{enumi}{1}
\tightlist
\item
  The Lebesgue measure on \(\mathbf{R}\) is not bounded; indeed, for every integer \(n > 0\) there exists a function \(f \in \mathcal{K}(\mathbf{R}; \mathbf{C})\) with values in {[}0,1{]} and equal to 1 on the interval \([-n, n]\) (No.~2, Lemma 1); thus \(\| f \| = 1\) and
\end{enumerate}

\[
\int_ {- \infty} ^ {+ \infty} f (x) d x \geqslant \int_ {- n} ^ {n} f (x) d x = 2 n,
\]

which proves that there does not exist any finite number \(M\) satisfying the relation (19).

\begin{enumerate}
\def\labelenumi{\arabic{enumi})}
\setcounter{enumi}{2}
\tightlist
\item
  On the real line \(\mathbf{R}\) , the mapping
\end{enumerate}

\[
f \mapsto \int_ {- \infty} ^ {+ \infty} \frac {f (x) d x}{1 + x ^ {2}}
\]

is a bounded measure since, for every function \(f \in \mathcal{K}(\mathbf{R};\mathbf{C})\) ,

\[
\left| \int_ {- \infty} ^ {+ \infty} \frac {f (x) d x}{1 + x ^ {2}} \right| \leqslant \| f \| \int_ {- \infty} ^ {+ \infty} \frac {d x}{1 + x ^ {2}} = \pi \cdot \| f \|.
\]

Since the relations \(\| f\| \leqslant 1\) and \(\| \overline{f}\| \leqslant 1\) are equivalent, it follows from (20) that

\[
\| \overline {{\mu}} \| = \| \mu \|\tag{21}
\]

for every measure \(\mu\) on X.

PROPOSITION 10. --- For every measure \(\mu\) on X,

\[
\| \mu \| = \sup _ {0 \leqslant f \leqslant 1, f \in \mathcal {K} (\mathrm{X}; \mathbf {R})} | \mu | (f).\tag{22}
\]

For, taking into account the formula (12) that defines the absolute value of a measure, the second member of (22) may be written

\[
\sup _ {0 \leqslant f \leqslant 1, f \in \mathcal {K} (\mathrm{X}; \mathbf {R})} \left(\sup _ {| g | \leqslant f, g \in \mathcal {K} (\mathrm{X}; \mathbf {C})} | \mu (g) |\right) = \sup _ {\| g \| \leqslant 1, g \in \mathcal {K} (\mathrm{X}; \mathbf {C})} | \mu (g) |.
\]

COROLLARY 1. --- For every measure \(\mu\) on X, the norms of \(\mu\) and \(|\mu|\) are equal; \(\mu\) is bounded if and only if \(|\mu|\) is bounded.

COROLLARY 2. --- For every measure \(\mu\) on a compact space \(X\) ,

\[
\| \mu \| = | \mu | (1) = \int d | \mu |.\tag{23}
\]

This formula will be generalized in Ch. IV, §4, No.~7.

On a compact space X, for every (complex) measure \(\mu\) on X the complex number \(\mu(1)\) is called the total mass of \(\mu\) . When \(\mu\) is positive, its total mass is thus equal to its norm. When \(\mu\) is a positive measure on a compact space X, of total mass equal to 1, one also says that its value \(\mu(f)\) for a continuous function \(f \in \mathcal{C}(\mathrm{X};\mathbf{C})\) is the mean of f with respect to the measure \(\mu\) .

COROLLARY 3. --- For every real measure \(\mu\) on a locally compact space \(X\) ,

\[
\| \mu \| = \sup _ {\| f \| \leqslant 1, f \in \mathcal{K} (\mathrm{X}; \mathbf {R})} | \mu (f) |.\tag{24}
\]

It suffices to make use of the formula (22) and the expression (9) for \(|\mu|(f)\) when \(\mu\) is a real measure and \(f \in \mathcal{K}_{+}(\mathrm{X})\) .

The set of bounded real measures is therefore the dual of the normed space \(\mathcal{K}(\mathrm{X};\mathbf{R})\) ; it is denoted \(\mathcal{M}^{1}(\mathrm{X},\mathbf{R})\) , or \(\mathcal{M}^{1}(\mathrm{X})\) if no confusion can result. The canonical injection \(\mathcal{M}^{1}(\mathrm{X},\mathbf{R})\to\mathcal{M}^{1}(\mathrm{X};\mathbf{C})\) is an isometry by virtue of (24).

PROPOSITION 11. --- If \(\mu\) and \(\nu\) are two positive measures on X, then \(\| \mu + \nu \| = \| \mu \| + \| \nu \|\) .

For, the functions \(f \in \mathcal{K}(\mathrm{X}; \mathbf{R})\) such that \(0 \leqslant f \leqslant 1\) form a directed set S for the relation \(\leqslant\) . For a positive measure \(\mu\) on X, it therefore follows from (22) and the monotone limit theorem that \(\| \mu \| = \lim_{f \in S} \mu(f)\) ; the conclusion of the proposition then follows at once.

COROLLARY 1. --- If \(\mu\) and \(\nu\) are two positive measures on X such that \(\mu \leqslant \nu\) , then \(\| \mu \| \leqslant \| \nu \|\) ; in particular, if \(\nu\) is bounded then so is \(\mu\) . Indeed, \(\| \nu \| = \| \mu \| + \| \nu - \mu \|\) .

COROLLARY 2. --- For every real measure \(\mu\) on X,

\[
\| \mu \| = \| \mu^ {+} \| + \| \mu^ {-} \|.
\]

For (Cor. 1 of Prop. 10), the norm of \(\mu\) is equal to that of \(|\mu| = \mu^{+} + \mu^{-}\) .

PROPOSITION 12. --- If \(\mu\) is a bounded measure on X and if \(g\) is a bounded continuous mapping of X into C, then the measure \(g \cdot \mu\) is bounded and \(\| g \cdot \mu \| \leqslant \| g \| \cdot \| \mu \|\) .

For every function \(f\in \mathcal{K}(\mathrm{X};\mathbf{C})\)

\[
| \mu (f g) | \leqslant \| \mu \| \cdot \| f g \| \leqslant \| \mu \| \cdot \| g \| \cdot \| f \|.
\]

